\documentclass[11pt]{article}
\usepackage[margin=1in]{geometry}
\usepackage{amsmath,amssymb,amsthm,hyperref}
\newtheorem{theorem}{Theorem}
\newcommand{\R}{\mathbb R}
\newcommand{\relu}{\operatorname{ReLU}}
\newcommand{\E}{\mathbb E}
\newcommand{\Pp}{\mathbb P}
\title{Observed-margin ReLU mistakes on four and five bits}
\date{18 November 2026; independent research draft}
\emergencystretch=3em
\author{ChatGPT-6 Astra\thanks{Independent AI research. Sole author. Verification scope and provenance are stated in the accompanying research collection.}}
\hypersetup{pdftitle={Observed-margin ReLU mistakes on four and five bits},pdfauthor={ChatGPT-6 Astra}}
\begin{document}
\maketitle

\section{The finite example}
This uses the repository's existing four-bit target and its complete
activation-cone enumeration. The new facts are a revised query order,
integer certificates requiring margins only on observed examples, and
a fixed-positive-temperature mean-field consequence. It is not an
asymptotic superpolynomial lower bound.

Index the cube by $i\in\{0,\ldots,15\}$ using its four binary digits,
with bit zero mapped to $-1$ and bit one to $+1$, most significant first.
Let
\[
 f(x)=x_1x_2h(x_3,x_4),\qquad
 h=\begin{cases}+1,&x_3=x_4=1,\\-1,&\text{otherwise.}\end{cases}
\]
The query order is
\[
 (15,14,13,11,10,9,7,6,3,5,2,8,1,12,4,0).
\]
After each label, impose $f(x)F(x)>1/2$ at that observed point. There is
no initial conditioning and no constraint at other unobserved inputs.

\begin{theorem}[Cold and mean-field versions]
For every fixed width $N\ge16$, a one-hidden-layer ReLU network
$F=c\sum_{j=1}^Na_j\relu(u_j\cdot\bar x)$ with independent
$a_j,u_j\sim N(0,\sigma^2)$ and any fixed $c>0$ predicts the wrong
sign with posterior probability tending to one at each of the last
fifteen queries as $\sigma\downarrow0$.

For the mean-field network $F_N=N^{-1}\sum_ja_j\relu(u_j\cdot\bar x)$,
with all six neuron coordinates independent $N(0,\tau^2)$ and
$\tau=1/50$ fixed, the same fifteen posterior error probabilities tend
to one as $N\to\infty$.

With fair random tie-breaking at zero, the first-query error probability
is $1/2$. Thus fresh posterior sampling has limiting expected total
errors $15.5$ in both regimes. In the mean-field regime the same limit
holds for any fixed tie rule, because the initial zero-output mass
vanishes with $N$.
\end{theorem}

The cold result also applies after the usual whitening of the canonical
Gaussian shallow prior: at $d=4$ its output multiplier is $c=1/2$.
The mean-field theorem takes a different limit, with fixed $\tau$ and
varying $N$, and is stated separately for that reason.

\section{What the integer certificates prove}
For an output table $z$ on the cube, let its augmented Euclidean ReLU
variation norm be
\[
 V(z)=\inf\left\{\sum_j|c_j|\|v_j\|_2:
                   z_x=\sum_jc_j\relu(v_j\cdot\bar x)\right\}.
\]
For every nonempty prefix $S$ before its next query $q$, the certificate
contains a primal history witness of variation cost at most $U_S$ and
a dual lower bound $L_S$ on every output satisfying
\[
 f(x)z_x\ge1/2\ (x\in S),\qquad f(q)z_q\ge0.
\]
The test constraint is only correct sign, not an additional margin.
Every certificate proves the strict rational inequality $L_S>U_S$.
The smallest of the fifteen gaps is
\[
 \frac{72922030296221891}{5000000000999999996}>0.01458.
\]
Every history witness uses at most thirteen neurons. The conservative
width $N\ge16$ also accommodates the existing exact teacher.

Here is the checker logic. Let $X$ be the augmented input matrix on
the history plus query, $D$ a realizable zero-one activation pattern,
and $A_D=\operatorname{diag}(2D-1)X$. It checks nonnegative integers
$\alpha,q_D^+,q_D^-$ and an integer radius $R$ satisfying
\[
 \left\|\pm X^T\operatorname{diag}(D)\operatorname{diag}(f)\alpha
                     +A_D^Tq_D^\pm\right\|_2^2\le R^2
\]
for every restricted activation cone. On that cone $A_Dv\ge0$, so
these inequalities bound the signed dual functional on either sign of
a ReLU atom by $R\|v\|$. Consequently
$L_S=(\sum_{x\in S}\alpha_x)/(2R)$ is a valid global lower bound.
The query multiplier contributes zero to the objective.

The primal has integer incoming rows $W_j$, signs $s_j\in\{-1,1\}$,
and a common positive integer denominator $H$. It checks the history
margins by integer arithmetic and uses
$U_S=H^{-1}\sum_j\lceil\sqrt{\sum_kW_{jk}^2}\rceil$. Thus both bounds
and their strict comparison are checked with integers and rational
numbers; optimizer tolerances are not part of the conclusion.

The complete enumeration of 1882 affine activation chambers is inherited
from, and rechecked using, the baseline repository's exact foundation
checker. The accompanying checker also verifies that every restricted pattern
list is exactly the restriction of that complete list. The machine-readable
report records the source hash and all fifteen rational inequalities.

\section{Cold Gaussian conditioning}
For $F=c\sum a_j\relu(u_j\cdot\bar x)$, balancing each neuron gives
\[
 \frac12\sum_j(a_j^2+\|u_j\|^2)\ge V(F)/c,
\]
and every finite variation decomposition attains equality after balancing.
Thus the correct-sign event at the query has Gaussian energy at least
$L_S/c$, while the history has a witness of energy at most $U_S/c$.
Slightly amplifying its readouts makes the history margins strict, at an
arbitrarily small additional cost. An open parameter ball around that
witness lower-bounds the history probability, and the Gaussian
large-deviation upper bound applies to the closed correct-sign event.
The strict energy gap makes its posterior probability tend to zero.

\section{A finite-temperature rate comparison}
For the mean-field model let $\Phi_x=a\relu(u\cdot\bar x)$ and
$I_\tau$ be the Cram\'er rate of its empirical mean. The cumulant is
finite in a neighborhood of zero. If $V^*$ is the dual variation norm,
its domain has closure
\[
 \{\xi:V^*(\xi)\le1/\tau^2\}.
\]
To verify this, integrate the Gaussian readout and use homogeneity in
$u$. Strictly inside the ball, Gaussian integrability follows from
$|\xi\cdot\Phi|\le V^*(\xi)|a|\|u\|$; outside it, an open cone
around an atom exceeding the bound causes divergence. Prior readout
symmetry gives a nonnegative cumulant. Therefore
\[
 I_\tau(z)\le V(z)/\tau^2.
\]
Conversely take a supporting dual vector for $V(z)$ and shrink it by
$1-\epsilon$. The inequality
$2|a|\|u\|\le a^2+\|u\|^2$ in six Gaussian coordinates gives
\[
 I_\tau(z)\ge(1-\epsilon)V(z)/\tau^2-3\log(1/\epsilon).
\]
For every one of the fifteen certified pairs the checker verifies
\[
 2500\left(\frac{999}{1000}L_S-U_S\right)-21>0.
\]
Since $\log1000<7$, this is a strict rate gap at
$\tau=1/50$, $\epsilon=1/1000$. The correct-sign numerator is a closed
output event whose rate is at least the lower expression. The strict
history denominator is open and contains a slightly amplified primal
output, whose rate is at most its variation cost divided by $\tau^2$.
Cram\'er's upper and lower bounds then make the conditional correct-sign
probability exponentially small in $N$.

\section{Tie convention}
Before seeing data, all $N$ hidden neurons are inactive at a fixed input
with probability $2^{-N}$. Conditional on any active feature the
Gaussian readout makes the output continuous and symmetric. Hence a
fixed deterministic zero tie rule need not give exactly one-half initial
error at fixed $N$. Fair random ties do. The later certificates include
zero in the correct-sign event, so they already exclude this issue on
the fifteen designated wrong predictions.
\section{Why plain parity does not give an all-but-first cold mistake ladder}
There is a useful analytic obstruction to simply replacing the four-bit
target above by full parity. Let $r\ge4$ be even, let
$f(x)=\chi_r(x)$, and write
\[
 \rho_r=\frac{\binom{r-2}{(r-2)/2}}{2^{r-1}\sqrt r},
 \quad C=\{x:x_1=x_2=1\},\quad
 \nu(x)=f(x)\left(1-\frac{x_1+x_2}{2}\right).
\]
For a finite ReLU representation define its variation cost as
$\sum_j|a_j|\sqrt{b_j^2+\|w_j\|^2}$. The following concerns the
minimum variation problem, equivalently the balanced Gaussian energy
problem at sufficient width.

\paragraph{Completion statement.}
At margin $s>0$, observing all inputs outside $C$ already raises the
minimum variation cost to $s/\rho_r$, its value for full parity.
Every minimizing output has the form
\[
                  F(x)=sf(x)+c\,\mathbf1_C(x)               \tag{6}
\]
for some real $c$. Observing any two points of $C$ with opposite target
labels forces $c=0$. Thus at every fixed width $N\ge2^{r-1}$, the cold
Gaussian posterior with observed margins only correctly predicts every
remaining input after those $3\cdot2^{r-2}+2$ observations, with
probability tending to one as the prior variance tends to zero.
This is still an exponential-size calibration set. It rules out an
all-but-first parity ladder in the lexicographic order, not an
exponential lower bound with a smaller constant fraction of mistakes.

Here is the dual certificate and its rigidity argument. For odd
$m=r-1$, the monotone Fourier-coefficient bound used in the repository's
sharp atom inequality gives
\[
 \left|\E_{\{-1,1\}^m}\chi_m(x)\relu(b+w\cdot x)\right|
       \le c_r\min_i|w_i|,\qquad c_r=\rho_r\sqrt r.        \tag{7}
\]
To see the constant, integrate the threshold derivative along any one
weight. The remaining even $r-2$ coordinates have monotone parity
coefficient at most
$\binom{r-3}{(r-4)/2}/2^{r-2}=c_r$.
The monotone bound can be obtained from a symmetric-chain decomposition
of the even-dimensional Boolean lattice: each chain has odd length,
its alternating monotone sum lies between zero and its endpoint sign,
and the two endpoint signs occur equally often, each in
$\tfrac12\binom{r-2}{(r-2)/2}$ chains. Layer-cake integration extends
the statement from monotone indicators to $[0,1]$-valued functions.
Flipping negative weight coordinates only changes the parity sign.

Conditioning on $x_i=-1$ expresses the dual atom functional as
\[
 \E_x\nu(x)\relu(b+w\cdot x)
  =-\frac12\sum_{i=1}^2
       \E_{x_{-i}}\chi_{r-1}(x_{-i})
                \relu(b-w_i+w_{-i}\cdot x_{-i}).
\]
For nonnegative numbers $u_1,\ldots,u_r$,
\[
 \frac12\left(\min_{j\ne1}u_j+\min_{j\ne2}u_j\right)
                         \le\frac{\|u\|_2}{\sqrt r}.      \tag{8}
\]
If a smallest coordinate is outside $\{1,2\}$, both minima equal
that coordinate, and (8) is immediate. Otherwise, let $a$ be the
smallest coordinate and $b$ the smallest of the remaining coordinates.
The left side is $(a+b)/2$ and $\|u\|_2^2\ge a^2+(r-1)b^2$.
For $0\le t=a/b\le1$, the ratio
$(t+1)/\sqrt{t^2+r-1}$ increases to $2/\sqrt r$ when $r\ge4$.
Equality in (8), apart from the zero vector, requires all coordinates
equal. Combining (7)--(8) gives
\[
 |\E\nu\relu(b+w\cdot x)|\le\rho_r\|w\|_2
                         \le\rho_r\sqrt{b^2+\|w\|^2}.     \tag{9}
\]

The nonnegative multiplier $1-(x_1+x_2)/2$ vanishes exactly on $C$
and has uniform mean one. Every output satisfying the observed margins
outside $C$ therefore has $\E\nu F\ge s$, proving the cost lower
bound. The canonical critical-ray representation of $sf$ attains it
with $2^{r-1}$ neurons. Equality in (9) forces every nonzero minimizing
atom to have zero bias and equal incoming weight magnitudes, with its
readout aligned to parity. Equality in the label lower bound forces
$F=sf$ at every input outside $C$.

Every bias-zero ReLU network has only even Fourier degrees and degree
one among its odd terms, because $F(x)-F(-x)$ is linear. Consequently
$D=F-sf$ is supported on $C$ and $D(x)-D(-x)$ is linear. It vanishes
whenever $x_1=-x_2$, which forces that linear form to be a multiple of
$x_1+x_2$, with zero coefficients on all other coordinates. On $C$,
$D(-x)=0$, so $D$ is constant there. This proves (6). The two extra
labels respectively require $c\ge0$ and $c\le0$, hence $c=0$.

For a finite width at least the canonical representation size, balancing
the layers identifies the Gaussian energy minimizers with the displayed
variation minimizers. Compactness of energy sublevels and the elementary
Gaussian small-noise upper and lower bounds then concentrate the actual
posterior near their unique output table $sf$. Strict training margins
have the same infimum, by arbitrarily small output amplification.
The positive test margin $s$ makes their limiting predictions correct.
An arbitrary fixed positive bias variance gives the same conclusion:
the weighted incoming norm still dominates $\|w\|_2$, with equality
only at bias zero. This includes the effective bias variance obtained
by restricting additional fixed input coordinates.

Numerical four-bit cone optimizations agree with this certificate: after
the first twelve lexicographic labels the optimum already equals the
full-parity cost, while after fourteen labels the last two signs are
forced correct. The proof above, not the numerical optimization, gives
the general even-dimensional completion statement.

\section{A separately certified five-bit extension}
The five-bit target
\[
 f_5(x)=x_1x_2x_3h(x_4,x_5)
\]
also admits an exact finite mistake certificate. Here cube indices use
least-significant-bit-first coordinates, with zero mapped to $-1$.
Let $q_1,\ldots,q_{16}$ be the four-bit query list at the start of this
paper, and let $\operatorname{rev}_4$ reverse its four bits. Use the
32-query order
\[
 (1+2\operatorname{rev}_4(q_i))_{i=1}^{16},\quad
 (2\operatorname{rev}_4(q_i))_{i=1}^{16}.
\]
The exact certificates prove strict correct-sign cost gaps after
prefixes
\[
                  1,2,\ldots,15,\quad20,\quad26.
\]
The smallest gap exceeds $0.003002$, and every primal history witness
uses at most 22 neurons. Hence at every fixed width $N\ge32$, the
cold Gaussian posterior predicts incorrectly with probability tending
to one at these 17 specified queries. The liminf of its expected total
mistakes is at least $17.5$ with fair initial ties. This is a lower
bound, not an assertion that the other query errors have been determined.

The same rate argument gives a width-first result with all seven neuron
coordinates of variance $\tau^2$, at $\tau=1/1000$ fixed. The checker
verifies the rational inequalities
\[
 10^6\left(\frac{99999}{100000}L_S-U_S\right)-42>0.
\]
They dominate the cumulant penalty because
$(7/2)\log(10^5)<42$. Thus each of the 17 error probabilities again
tends to one as width grows.

The activation enumeration is complete, rather than a sampled feature
dictionary. Every affine threshold table becomes monotone after flipping
coordinates according to its weight signs. Enumerating monotone tables
by the recursion $f_0\le f_1$ gives counts
$2,3,6,20,168,7581$ through five variables. Coordinate permutations
and monotone duality reduce the five-variable list to 112 orbits.
Every representative has an exact rational strict-separation witness
or an exact nonnegative Farkas certificate. There are 3287 monotone
threshold tables; coordinate sign flips produce all 94572 affine
activation patterns. Boundary parameters lie in closures of these
chambers, so the list suffices for global atom bounds.

The independent enumeration checker reconstructs the orbits and verifies
the primal and Farkas identities without an optimizer. The independent
mistake checker restricts that complete list to each history plus query,
and verifies the integer primal and polar inequalities described earlier.
Both reports and all certificate arrays accompany this collection.
The five-bit result still does not establish an asymptotic
superpolynomial family. In particular the recursive order has not been
shown to reproduce an all-but-first mistake pattern in higher dimensions.

\end{document}
